\documentclass[11pt,a4paper]{article}
% preamble.tex — estilo común de todos los apuntes Froth.
% Cada apunte hace % preamble.tex — estilo común de todos los apuntes Froth.
% Cada apunte hace % preamble.tex — estilo común de todos los apuntes Froth.
% Cada apunte hace \input{preamble} justo después de \documentclass.
\usepackage[margin=2.4cm]{geometry}
\usepackage{xcolor}
\definecolor{litblue}{RGB}{31,79,121}
\definecolor{litgreen}{RGB}{27,94,32}
\definecolor{litorange}{RGB}{230,81,0}
\usepackage{parskip}
\usepackage{titlesec}
\titleformat{\section}{\Large\bfseries\color{litblue}}{\thesection}{0.6em}{}
\titleformat{\subsection}{\large\bfseries\color{litblue}}{\thesubsection}{0.6em}{}
\usepackage{tcolorbox}
\usepackage{fancyhdr}
\pagestyle{fancy}
\fancyhf{}
\fancyhead[L]{\small\color{litblue}Froth — Apuntes de ML}
\fancyhead[R]{\small\thepage}
\renewcommand{\headrulewidth}{0.4pt}
\usepackage[colorlinks=true,linkcolor=litblue,urlcolor=litblue]{hyperref}

% Cabecera de cada apunte: \cabecera{numero}{titulo}
\newcommand{\cabecera}[2]{%
  \begin{tcolorbox}[colback=litblue,coltext=white,colframe=litblue,arc=2mm]
    {\Large\bfseries Apunte #1 — #2}\\[3pt]
    {\small Proyecto Froth · aprender Machine Learning construyendo}
  \end{tcolorbox}\vspace{0.8em}}

% Cajas temáticas
\newtcolorbox{concepto}[1]{colback=blue!4,colframe=litblue,title={Concepto: #1},arc=1.5mm}
\newtcolorbox{practica}{colback=green!5,colframe=litgreen,title={Pruébalo tú (teclado en mano)},arc=1.5mm}
\newtcolorbox{ojo}{colback=orange!8,colframe=litorange,title={Ojo / error típico},arc=1.5mm}
\newtcolorbox{resumen}{colback=gray!8,colframe=gray!60,title={En una frase},arc=1.5mm}

\newcommand{\codigo}[1]{\texttt{#1}}
 justo después de \documentclass.
\usepackage[margin=2.4cm]{geometry}
\usepackage{xcolor}
\definecolor{litblue}{RGB}{31,79,121}
\definecolor{litgreen}{RGB}{27,94,32}
\definecolor{litorange}{RGB}{230,81,0}
\usepackage{parskip}
\usepackage{titlesec}
\titleformat{\section}{\Large\bfseries\color{litblue}}{\thesection}{0.6em}{}
\titleformat{\subsection}{\large\bfseries\color{litblue}}{\thesubsection}{0.6em}{}
\usepackage{tcolorbox}
\usepackage{fancyhdr}
\pagestyle{fancy}
\fancyhf{}
\fancyhead[L]{\small\color{litblue}Froth — Apuntes de ML}
\fancyhead[R]{\small\thepage}
\renewcommand{\headrulewidth}{0.4pt}
\usepackage[colorlinks=true,linkcolor=litblue,urlcolor=litblue]{hyperref}

% Cabecera de cada apunte: \cabecera{numero}{titulo}
\newcommand{\cabecera}[2]{%
  \begin{tcolorbox}[colback=litblue,coltext=white,colframe=litblue,arc=2mm]
    {\Large\bfseries Apunte #1 — #2}\\[3pt]
    {\small Proyecto Froth · aprender Machine Learning construyendo}
  \end{tcolorbox}\vspace{0.8em}}

% Cajas temáticas
\newtcolorbox{concepto}[1]{colback=blue!4,colframe=litblue,title={Concepto: #1},arc=1.5mm}
\newtcolorbox{practica}{colback=green!5,colframe=litgreen,title={Pruébalo tú (teclado en mano)},arc=1.5mm}
\newtcolorbox{ojo}{colback=orange!8,colframe=litorange,title={Ojo / error típico},arc=1.5mm}
\newtcolorbox{resumen}{colback=gray!8,colframe=gray!60,title={En una frase},arc=1.5mm}

\newcommand{\codigo}[1]{\texttt{#1}}
 justo después de \documentclass.
\usepackage[margin=2.4cm]{geometry}
\usepackage{xcolor}
\definecolor{litblue}{RGB}{31,79,121}
\definecolor{litgreen}{RGB}{27,94,32}
\definecolor{litorange}{RGB}{230,81,0}
\usepackage{parskip}
\usepackage{titlesec}
\titleformat{\section}{\Large\bfseries\color{litblue}}{\thesection}{0.6em}{}
\titleformat{\subsection}{\large\bfseries\color{litblue}}{\thesubsection}{0.6em}{}
\usepackage{tcolorbox}
\usepackage{fancyhdr}
\pagestyle{fancy}
\fancyhf{}
\fancyhead[L]{\small\color{litblue}Froth — Apuntes de ML}
\fancyhead[R]{\small\thepage}
\renewcommand{\headrulewidth}{0.4pt}
\usepackage[colorlinks=true,linkcolor=litblue,urlcolor=litblue]{hyperref}

% Cabecera de cada apunte: \cabecera{numero}{titulo}
\newcommand{\cabecera}[2]{%
  \begin{tcolorbox}[colback=litblue,coltext=white,colframe=litblue,arc=2mm]
    {\Large\bfseries Apunte #1 — #2}\\[3pt]
    {\small Proyecto Froth · aprender Machine Learning construyendo}
  \end{tcolorbox}\vspace{0.8em}}

% Cajas temáticas
\newtcolorbox{concepto}[1]{colback=blue!4,colframe=litblue,title={Concepto: #1},arc=1.5mm}
\newtcolorbox{practica}{colback=green!5,colframe=litgreen,title={Pruébalo tú (teclado en mano)},arc=1.5mm}
\newtcolorbox{ojo}{colback=orange!8,colframe=litorange,title={Ojo / error típico},arc=1.5mm}
\newtcolorbox{resumen}{colback=gray!8,colframe=gray!60,title={En una frase},arc=1.5mm}

\newcommand{\codigo}[1]{\texttt{#1}}

\begin{document}
\cabecera{07}{Vectorizar los 126 papers: la matriz de embeddings}

\section{De uno a todos}

En el Apunte 06 convertimos \emph{un} paper en su vector. Aquí convertimos los
\textbf{126 de golpe} y \textbf{guardamos} el resultado en disco. Esta es la diferencia
clave: antes solo mirábamos; ahora producimos un dato que las siguientes fases reutilizarán.

\begin{concepto}{matriz (N filas $\times$ M columnas)}
Un vector es una fila de números. Si apilas muchos vectores, obtienes una \textbf{matriz}:
una tabla de números. La nuestra tiene forma \codigo{(126, 768)}:
\begin{itemize}
  \item \textbf{126 filas} --- una por paper.
  \item \textbf{768 columnas} --- las coordenadas de cada paper.
\end{itemize}
La \emph{fila i} de la matriz es el vector del \emph{paper i} de la tabla: mismo orden. Esa
alineación es sagrada --- así sabemos qué vector pertenece a qué paper.
\end{concepto}

\section{Procesar en lote (batch): por qué es más rápido}

El código no llama al modelo 126 veces (una por paper). Le pasa la \textbf{lista completa}
de textos y el modelo los procesa en \emph{lotes} (viste la barra \codigo{Batches: 4/4}).

\begin{concepto}{batch (lote)}
Procesar en lote significa dar muchos ejemplos juntos para que el modelo aproveche mejor el
procesador (hace las cuentas en paralelo). Uno por uno funcionaría, pero sería más lento.
Nuestros 126 papers tardaron $\sim$45 segundos en CPU; con una tarjeta gráfica sería aún
más rápido, pero para este tamaño no hace falta.
\end{concepto}

\section{Guardar en disco: el archivo \codigo{.npy}}

Vectorizar cuesta segundos; no queremos repetirlo cada vez. Por eso guardamos la matriz en
\codigo{2\_Datos/embeddings/vectors.npy}.

\begin{concepto}{formato .npy (NumPy)}
\codigo{.npy} es el formato propio de NumPy para guardar arrays/matrices de números tal
cual, sin convertir a texto. Se escribe con \codigo{np.save()} y se relee al instante con
\codigo{np.load()}. Nuestra matriz de 126$\times$768 pesa apenas $\sim$378 KB. Las fases
siguientes (reducir a 2D, agrupar) leerán este archivo en vez de recalcular.
\end{concepto}

\begin{ojo}
\textbf{Regla de oro del pipeline:} cada fase deja su resultado guardado y la siguiente lo
lee. Papers $\rightarrow$ \codigo{papers.parquet}; vectores $\rightarrow$ \codigo{vectors.npy}.
Si algo se rompe más adelante, no tienes que volver a bajar papers ni recalcular embeddings:
retomas desde el último archivo guardado.
\end{ojo}

\section{Lo que quedó construido}

\begin{itemize}
  \item La función \codigo{embed\_papers(df)} en \codigo{froth/embed.py} (ya no está vacía).
  \item El archivo \codigo{2\_Datos/embeddings/vectors.npy} con la matriz 126$\times$768.
  \item Utilidades \codigo{save\_embeddings()} y \codigo{load\_embeddings()} para las
        fases que vienen.
\end{itemize}

\begin{resumen}
Apilamos los 126 vectores en una matriz (126$\times$768), la calculamos en lote (rápido) y la
guardamos en \codigo{vectors.npy} para no recalcular; la fila i siempre es el paper i.
\end{resumen}

\end{document}
